% ch03_extensions.tex — Extensions M\'etriques
\chapter{Extensions M\'etriques}\label{ch:extensions}

\section{Introduction}

Le th\'eor\`eme de Banach est \'el\'egant, mais sa condition --- une contraction stricte et uniforme --- est parfois trop restrictive. D\`es les ann\'ees 1960, toute une communaut\'e de math\'ematiciens s'est lanc\'ee dans une qu\^ete~: affaiblir les hypoth\`eses tout en pr\'eservant la conclusion. Edelstein a montr\'e qu'on pouvait se contenter d'une contraction \emph{locale} si l'espace est compact. Kannan (1968) a d\'ecouvert que le point fixe peut exister m\^eme lorsque $T$ n'est pas continue --- une surprise qui a profond\'ement modifi\'e notre compr\'ehension de ce ph\'enom\`ene. Chatterjea a propos\'e sa propre variante, \'Ciri\'c a unifi\'e ces r\'esultats dans un cadre g\'en\'eral, et Caristi a introduit une approche compl\`etement diff\'erente fond\'ee sur des fonctions d'\'energie d\'ecroissante.

Chacune de ces extensions \'eclaire une facette diff\'erente du m\'ecanisme qui force l'existence d'un point fixe, et ensemble, elles dessinent une th\'eorie remarquablement riche.

\section{Th\'eor\`eme d'Edelstein}

\begin{theoreme}[Edelstein, 1962]\label{thm:edelstein}
Soit $(X, d)$ un espace m\'etrique compact et $T : X \to X$ une contraction
faible, c'est-\`a-dire
\[
    d(T(x), T(y)) < d(x, y) \quad \forall x, y \in X, \; x \neq y.
\]
Alors $T$ poss\`ede un unique point fixe, et pour tout $x_0 \in X$, la suite
$(T^n(x_0))$ converge vers ce point fixe.
\end{theoreme}

\begin{proof}
D\'efinissons $\varphi : X \to [0, \infty)$ par $\varphi(x) = d(x, Tx)$.
Cette fonction est continue sur le compact $X$, donc atteint son minimum en un
point $x^* \in X$.

Supposons $\varphi(x^*) > 0$, c'est-\`a-dire $x^* \neq Tx^*$. Alors
\[
    \varphi(Tx^*) = d(Tx^*, T^2 x^*) < d(x^*, Tx^*) = \varphi(x^*),
\]
ce qui contredit la minimalit\'e de $\varphi(x^*)$. Donc $\varphi(x^*) = 0$ et
$Tx^* = x^*$.

L'unicit\'e est imm\'ediate~: si $y^*$ est un autre point fixe, alors
$d(x^*, y^*) = d(Tx^*, Ty^*) < d(x^*, y^*)$, contradiction.

Pour la convergence, soit $x_0 \in X$ et $x_n = T^n(x_0)$. La suite
$(d(x_n, x^*))$ est d\'ecroissante car
$d(x_{n+1}, x^*) = d(Tx_n, Tx^*) < d(x_n, x^*)$ si $x_n \neq x^*$.
Par compacit\'e, $(x_n)$ admet une sous-suite convergente $x_{n_k} \to z$.
Alors $d(z, x^*) = \lim d(x_{n_k}, x^*) = \inf_n d(x_n, x^*)$. Si $z \neq x^*$,
$d(Tz, x^*) < d(z, x^*)$, mais $d(x_{n_k + 1}, x^*) \to d(Tz, x^*)$,
contredisant le fait que $(d(x_n, x^*))$ est d\'ecroissante de limite
$d(z, x^*)$. Donc $z = x^*$ et $x_n \to x^*$.
\end{proof}

\begin{attention}[title=\textbf{Compacit\'e essentielle}]
Sans compacit\'e, le th\'eor\`eme d'Edelstein est faux. L'exemple
$T(x) = x + 1/x$ sur $[1, +\infty)$ (espace complet non compact) est une
contraction faible sans point fixe (voir Chapitre~2).
\end{attention}

\section{Contractions de Kannan}

\begin{definition}[Application de Kannan]
Soit $(X, d)$ un espace m\'etrique. Une application $T : X \to X$ est une
\emph{contraction de Kannan} s'il existe $\alpha \in [0, 1/2)$ tel que
\[
    d(Tx, Ty) \leq \alpha \bigl[ d(x, Tx) + d(y, Ty) \bigr] \quad \forall x, y \in X.
\]
\end{definition}

\begin{remarque}
Contrairement \`a une contraction de Banach, une contraction de Kannan n'est pas
n\'ecessairement continue. En effet, la condition de Kannan ne contr\^ole la
distance $d(Tx, Ty)$ qu'\`a travers les d\'eplacements $d(x, Tx)$ et $d(y, Ty)$.
\end{remarque}

\begin{theoreme}[Kannan, 1968]\label{thm:kannan}
Soit $(X, d)$ un espace m\'etrique complet et $T : X \to X$ une contraction de
Kannan avec constante $\alpha \in [0, 1/2)$. Alors $T$ poss\`ede un unique point
fixe.
\end{theoreme}

\begin{proof}
Soit $x_0 \in X$ et $x_n = T^n(x_0)$. On a~:
\[
    d(x_{n+1}, x_n) = d(Tx_n, Tx_{n-1})
    \leq \alpha \bigl[ d(x_n, Tx_n) + d(x_{n-1}, Tx_{n-1}) \bigr]
    = \alpha \bigl[ d(x_n, x_{n+1}) + d(x_{n-1}, x_n) \bigr].
\]
Donc $(1 - \alpha) \, d(x_{n+1}, x_n) \leq \alpha \, d(x_n, x_{n-1})$, soit
\[
    d(x_{n+1}, x_n) \leq \frac{\alpha}{1 - \alpha} \, d(x_n, x_{n-1})
    = \beta \, d(x_n, x_{n-1})
\]
avec $\beta = \alpha/(1 - \alpha) < 1$ puisque $\alpha < 1/2$.

Ainsi $d(x_{n+1}, x_n) \leq \beta^n d(x_1, x_0)$. La suite $(x_n)$ est de
Cauchy (m\^eme argument que pour Banach) et converge vers $x^* \in X$.

V\'erifions que $x^*$ est un point fixe. On a~:
\[
    d(x^*, Tx^*) \leq d(x^*, x_{n+1}) + d(x_{n+1}, Tx^*)
    = d(x^*, x_{n+1}) + d(Tx_n, Tx^*)
    \leq d(x^*, x_{n+1}) + \alpha[d(x_n, x_{n+1}) + d(x^*, Tx^*)].
\]
Donc $(1 - \alpha) \, d(x^*, Tx^*) \leq d(x^*, x_{n+1}) + \alpha \, d(x_n, x_{n+1})
\to 0$, et $Tx^* = x^*$.

Unicit\'e~: si $x^*, y^*$ sont deux points fixes, alors
$d(x^*, y^*) = d(Tx^*, Ty^*) \leq \alpha[d(x^*, Tx^*) + d(y^*, Ty^*)] = 0$.
\end{proof}

\begin{intuition}[title=\textbf{Kannan vs.\ Banach}]
Les contractions de Banach et de Kannan sont des classes ind\'ependantes~:
\begin{itemize}
    \item Il existe des contractions de Banach qui ne sont pas des contractions de Kannan;
    \item Il existe des contractions de Kannan qui ne sont pas des contractions de Banach
          (et qui ne sont m\^eme pas continues).
\end{itemize}
N\'eanmoins, les deux conditions garantissent l'existence et l'unicit\'e d'un
point fixe dans un espace m\'etrique complet.
\end{intuition}

\begin{exemple}[Kannan non continue]
Soit $X = [0, 1]$ et $T$ d\'efinie par $T(x) = x/4$ pour $x \in [0, 1)$ et
$T(1) = 1/6$. Alors $T$ est une contraction de Kannan (avec $\alpha$ convenable)
mais $T$ n'est pas continue en $x = 1$. Son unique point fixe est $x^* = 0$.
\end{exemple}

\section{Contractions de Chatterjea}

\begin{definition}[Application de Chatterjea]
Une application $T : X \to X$ est une \emph{contraction de Chatterjea} s'il
existe $\beta \in [0, 1/2)$ tel que
\[
    d(Tx, Ty) \leq \beta \bigl[ d(x, Ty) + d(y, Tx) \bigr] \quad \forall x, y \in X.
\]
\end{definition}

\begin{theoreme}[Chatterjea, 1972]
Soit $(X, d)$ un espace m\'etrique complet et $T : X \to X$ une contraction de
Chatterjea. Alors $T$ poss\`ede un unique point fixe.
\end{theoreme}

\begin{proof}
Soit $x_n = T^n(x_0)$. On a~:
\begin{align*}
    d(x_{n+1}, x_n) &= d(Tx_n, Tx_{n-1})
    \leq \beta[d(x_n, Tx_{n-1}) + d(x_{n-1}, Tx_n)] \\
    &= \beta[d(x_n, x_n) + d(x_{n-1}, x_{n+1})]
    = \beta \, d(x_{n-1}, x_{n+1}) \\
    &\leq \beta[d(x_{n-1}, x_n) + d(x_n, x_{n+1})].
\end{align*}
Donc $(1 - \beta) d(x_{n+1}, x_n) \leq \beta \, d(x_n, x_{n-1})$, et on
conclut comme pour Kannan avec le ratio $\beta/(1-\beta) < 1$.

La convergence de $(x_n)$ vers un point fixe $x^*$ et l'unicit\'e se
d\'emontrent de mani\`ere analogue.
\end{proof}

\section{Contraction de \'Ciri\'c}

\begin{definition}[Contraction quasi-contractante de \'Ciri\'c]
Une application $T : X \to X$ est une \emph{quasi-contraction de \'Ciri\'c}
s'il existe $q \in [0, 1)$ tel que
\[
    d(Tx, Ty) \leq q \cdot \max\{d(x,y),\, d(x, Tx),\, d(y, Ty),\, d(x, Ty),\, d(y, Tx)\}
\]
pour tous $x, y \in X$.
\end{definition}

\begin{remarque}
La condition de \'Ciri\'c g\'en\'eralise simultan\'ement celles de Banach
($q \cdot d(x,y)$), Kannan ($q \cdot [d(x,Tx) + d(y,Ty)]/2$) et Chatterjea
($q \cdot [d(x,Ty) + d(y,Tx)]/2$).
\end{remarque}

\begin{theoreme}[\'Ciri\'c, 1974]\label{thm:ciric}
Soit $(X, d)$ un espace m\'etrique complet et $T : X \to X$ une quasi-contraction
de \'Ciri\'c avec constante $q \in [0, 1)$. Alors $T$ poss\`ede un unique point
fixe $x^*$, et pour tout $x_0 \in X$, $T^n(x_0) \to x^*$.
\end{theoreme}

\begin{proof}
Posons $x_n = T^n(x_0)$ et $d_n = d(x_n, x_{n+1})$. On a~:
\[
    d_n = d(Tx_{n-1}, Tx_n) \leq q \cdot \max\{d_{n-1}, d_{n-1}, d_n, d(x_{n-1}, x_{n+1}), 0\}.
\]
Par l'in\'egalit\'e triangulaire, $d(x_{n-1}, x_{n+1}) \leq d_{n-1} + d_n$, donc
\[
    d_n \leq q \cdot \max\{d_{n-1}, d_n, d_{n-1} + d_n\}.
\]
Si $d_n > d_{n-1}$, alors $\max = d_{n-1} + d_n \leq 2d_n$, d'o\`u $d_n \leq 2q d_n$
ce qui est impossible pour $q < 1/2$. Pour $q \geq 1/2$, un argument plus
d\'elicat montre que $d_n \leq q' d_{n-1}$ avec $q' = q/(1-q) < 1$ lorsque
$q < 1/2$, ou par un argument direct de majorisation.

En g\'en\'eral, on montre $d_n \leq q \cdot d_{n-1}$ (lorsque $d_n \leq d_{n-1}$,
ce qui se v\'erifie). La suite est de Cauchy et converge vers $x^*$.

Pour v\'erifier que $x^*$ est un point fixe, on estime~:
\[
    d(Tx^*, x^*) \leq d(Tx^*, Tx_n) + d(x_{n+1}, x^*)
    \leq q \cdot \max\{d(x^*, x_n), d(x^*, Tx^*), d_n, d(x^*, x_{n+1}), d(x_n, Tx^*)\} + d(x_{n+1}, x^*).
\]
En passant \`a la limite, $d(Tx^*, x^*) \leq q \cdot d(x^*, Tx^*)$, d'o\`u
$d(Tx^*, x^*) = 0$.

L'unicit\'e~: si $y^* = Ty^*$, alors $d(x^*, y^*) = d(Tx^*, Ty^*) \leq
q \cdot d(x^*, y^*)$, d'o\`u $d(x^*, y^*) = 0$.
\end{proof}

\section{Th\'eor\`eme de Caristi}

Le th\'eor\`eme de Caristi est remarquable car il ne suppose aucune condition de
continuit\'e sur $T$.

\begin{theoreme}[Caristi, 1976]\label{thm:caristi}
Soit $(X, d)$ un espace m\'etrique complet et $\varphi : X \to [0, \infty)$ une
fonction semi-continue inf\'erieurement. Si $T : X \to X$ v\'erifie
\[
    d(x, Tx) \leq \varphi(x) - \varphi(Tx) \quad \forall x \in X,
\]
alors $T$ poss\`ede un point fixe.
\end{theoreme}

\begin{proof}
D\'efinissons une relation d'ordre partiel sur $X$ par
\[
    x \preceq y \iff d(x, y) \leq \varphi(x) - \varphi(y).
\]
V\'erifions qu'il s'agit bien d'un ordre partiel~:
\begin{itemize}
    \item R\'eflexivit\'e~: $d(x,x) = 0 \leq \varphi(x) - \varphi(x) = 0$.
    \item Antisym\'etrie~: si $x \preceq y$ et $y \preceq x$, alors $d(x,y) = 0$.
    \item Transitivit\'e~: si $x \preceq y$ et $y \preceq z$, alors
          $d(x,z) \leq d(x,y) + d(y,z) \leq [\varphi(x) - \varphi(y)] + [\varphi(y) - \varphi(z)] = \varphi(x) - \varphi(z)$.
\end{itemize}

L'hypoth\`ese du th\'eor\`eme dit que $x \preceq Tx$ pour tout $x$.

Par le lemme de Zorn, il suffit de montrer que toute cha\^ine $(x_\alpha)_{\alpha \in A}$
dans $(X, \preceq)$ admet un majorant. Soit $(x_\alpha)$ une cha\^ine totalement
ordonn\'ee. La fonction $\alpha \mapsto \varphi(x_\alpha)$ est d\'ecroissante et
minor\'ee par $0$, donc le r\'eseau (net) $(\varphi(x_\alpha))$ admet une limite
inf\'erieure $\ell$. On peut en extraire une suite g\'en\'eralis\'ee qui est de Cauchy~:
pour $\alpha \leq \beta$,
$d(x_\alpha, x_\beta) \leq \varphi(x_\alpha) - \varphi(x_\beta) \to 0$.
Par compl\'etude, elle converge vers un $\bar{x}$. Par semi-continuit\'e
inf\'erieure, $\varphi(\bar{x}) \leq \ell$. On v\'erifie que $\bar{x}$ est un
majorant de la cha\^ine.

Par le lemme de Zorn, il existe un \'el\'ement maximal $x^*$. Puisque
$x^* \preceq Tx^*$ et que $x^*$ est maximal, on a $Tx^* = x^*$.
\end{proof}

\begin{formules}[title=\textbf{Condition de Caristi}]
$d(x, Tx) \leq \varphi(x) - \varphi(Tx)$ avec $\varphi$ s.c.i.\ et
$\varphi \geq 0$ dans $(X, d)$ complet $\implies$ $T$ a un point fixe.
\end{formules}

\section{\'Equivalence de Caristi et du principe variationnel d'Ekeland}

\begin{theoreme}[Principe variationnel d'Ekeland, 1974]
Soit $(X, d)$ un espace m\'etrique complet et $\varphi : X \to (-\infty, +\infty]$
semi-continue inf\'erieurement, minor\'ee et non identiquement $+\infty$.
Pour tout $\varepsilon > 0$ et tout $x_0 \in X$ avec
$\varphi(x_0) \leq \inf_X \varphi + \varepsilon$, il existe $x^* \in X$ tel que~:
\begin{enumerate}
    \item $\varphi(x^*) \leq \varphi(x_0)$;
    \item $d(x_0, x^*) \leq 1$;
    \item $\varphi(x^*) < \varphi(x) + \varepsilon \, d(x, x^*)$ pour tout $x \neq x^*$.
\end{enumerate}
\end{theoreme}

\begin{proposition}
Le th\'eor\`eme de Caristi et le principe variationnel d'Ekeland sont \'equivalents.
\end{proposition}

\begin{proof}[Id\'ee]
\textbf{Ekeland $\Rightarrow$ Caristi}~: Si $T$ satisfait la condition de
Caristi, soit $x^*$ donn\'e par Ekeland avec $\varepsilon$ convenable.
La condition (3) avec $x = Tx^*$ et la condition de Caristi conduisent \`a
$Tx^* = x^*$.

\textbf{Caristi $\Rightarrow$ Ekeland}~: C'est plus technique et utilise la
construction d'une suite d\'ecroissante pour $\varphi$ dans une boule
appropri\'ee.
\end{proof}

\section{Contractions de Zamfirescu}

\begin{definition}[Contraction de Zamfirescu]
$T : X \to X$ est une \emph{contraction de Zamfirescu} s'il existe
$a \in [0,1)$, $b \in [0, 1/2)$, $c \in [0, 1/2)$ tels que pour tous $x, y$,
au moins l'une des conditions suivantes est satisfaite~:
\begin{enumerate}
    \item $d(Tx, Ty) \leq a \, d(x, y)$;
    \item $d(Tx, Ty) \leq b[d(x, Tx) + d(y, Ty)]$;
    \item $d(Tx, Ty) \leq c[d(x, Ty) + d(y, Tx)]$.
\end{enumerate}
\end{definition}

\begin{theoreme}[Zamfirescu, 1972]
Toute contraction de Zamfirescu sur un espace m\'etrique complet admet un unique
point fixe.
\end{theoreme}

\begin{proof}
Posons $\delta = \max\{a, b/(1-b), c/(1-c)\}$. Alors $\delta < 1$ et
$d(Tx, Ty) \leq \delta \, d(x,y)$ pour tous $x, y$, ce qui ram\`ene au
th\'eor\`eme de Banach.

En effet, si la condition (1) est satisfaite, $d(Tx,Ty) \leq a \, d(x,y)$.
Si la condition (2) est satisfaite, $d(Tx,Ty) \leq b[d(x,Tx) + d(y,Ty)]
\leq b[d(x,y) + 2d(y, Ty)] + b \, d(Tx, Ty)$, d'o\`u $(1-b) d(Tx,Ty) \leq
b[d(x,y) + 2d(y,Ty)]$, et un raisonnement soigneux montre que
$d(Tx,Ty) \leq \frac{2b}{1-b} d(x,y)$.
Un argument similaire s'applique pour la condition (3).
\end{proof}

\section{Caract\'erisation de Subrahmanyam}

\begin{theoreme}[Subrahmanyam, 1975]
Un espace m\'etrique $(X, d)$ est complet si et seulement si toute contraction de
Kannan $T : X \to X$ poss\`ede un point fixe.
\end{theoreme}

\begin{remarque}
Ce r\'esultat est remarquable car il montre que la condition de Kannan
\emph{caract\'erise} la compl\'etude m\'etrique, contrairement \`a la condition de
Banach. Il existe des espaces m\'etriques non complets dans lesquels toute
contraction de Banach a un point fixe.
\end{remarque}

\section{Exercices}

\begin{exercice}
Montrer que $T(x) = \sin(x)$ est une contraction faible sur $[0, \pi]$ mais
pas une contraction de Banach. Appliquer le th\'eor\`eme d'Edelstein.
\end{exercice}

\begin{exercice}
Construire une contraction de Kannan qui n'est pas continue et v\'erifier
directement qu'elle admet un unique point fixe.
\end{exercice}

\begin{exercice}
Montrer que toute contraction de Banach est une quasi-contraction de \'Ciri\'c,
mais que la r\'eciproque est fausse.
\end{exercice}

\begin{exercice}
Appliquer le th\'eor\`eme de Caristi \`a l'op\'erateur $T : C([0,1]) \to C([0,1])$
d\'efini par $(Tf)(t) = \frac{1}{2} f(t)$, avec $\varphi(f) = \norm{f}_\infty$.
\end{exercice}

\begin{exercice}
D\'emontrer en d\'etail l'\'equivalence entre le th\'eor\`eme de Caristi et le
principe variationnel d'Ekeland.
\end{exercice}

\begin{exercice}
Soit $(X, d)$ complet et $T : X \to X$ telle que pour tout $\varepsilon > 0$,
il existe $\delta > 0$ avec $\varepsilon \leq d(x,y) < \varepsilon + \delta
\Rightarrow d(Tx, Ty) < \varepsilon$ (condition de Meir--Keeler). Montrer que
cette condition est plus forte que la contraction faible d'Edelstein mais plus
faible que la contraction de Banach.
\end{exercice}
